% ---------------------------------------------------------------------------
% Appendix: Prompt Library and Intervention Inventory
% Sources: Engineering_simplicity/contents/appendix.tex (prompt sections +
%   tab:all-interventions), writeup/contents/appendix.tex:577-605 (menu +
%   clock-framing prompts), verified against the repo templates in
%   Engineering_simplicity/engineer_simplicity-main/rule_template/{auctions,DA}/
%   and rule_template/V10,V12/ (old labels), and against the pipeline logs in
%   engineer_simplicity-main/experiment_logs/ (run-status audit).
% Implements plan §7.8 (no phantom results) and plan §5.6/§5.7 (old->new
%   taxonomy mapping; intervention_proxy_breitmoser split).
% Prompt texts are quoted verbatim from the templates actually sent to the
%   models (including their typos); non-ASCII glyphs are transliterated.
% ---------------------------------------------------------------------------
\section{Prompt Library and Intervention Inventory}\label{app:prompts}

This appendix documents the prompt templates behind every design lever used in the paper, organized by the taxonomy of Section~\ref{sec:framework}: Family~A (extensive-form implementations), Family~B (mechanism descriptions: B1 menu restatement, B2 safety-exposing, B3 clock-framing), Family~C (reasoning scaffolds: C1 contingent, C2 forward planning, C3 beliefs), and Family~D (preference framings, Appendix~\ref{app:prospect}). We present the full base-rule prompts for both domains; for each intervention we show only the text that is \emph{added to or modified from} the corresponding base prompt, since the surrounding context (bidder identity and response format in auctions; preference ordering, priorities, popularity signal, and output format in DA) remains identical. Template variables (e.g., \texttt{\{\{student\_id\}\}}, \texttt{\{\{preference\_order\}\}}, \texttt{\{\{increment\}\}}) are populated at runtime. The simulation pipeline that embeds these prompts is described in Appendix~\ref{app:procedures}; the DA protocol details are in Appendix~\ref{app:da}. Section~\ref{app:prompts-inventory} closes with the complete inventory (Table~\ref{tab:intervention-inventory}), which maps every repo-level template name to its final taxonomy family and records whether its results appear in this paper.

\subsection{Base rules}\label{app:prompts-base}

\paragraph{SPSB auction (IPV baseline).} The canonical baseline for the auction ranking cells (repo template \path{private_second_price}):

\begin{figure}[H]
\begin{lstlisting}
In this game, you will participate in an auction for a prize
against {{num_bidders}} other bidders. You will play this game
for {{n}} rounds.
At the start of each round, bidders will see their value for
the prize, randomly drawn between $0 and ${{private}}, with
all values equally likely.
After learning your value, you will submit a bid privately at
the same time as the other bidders. Bids must be in
${{increment}} increments.
The highest bidder wins the prize and pays the second-highest
bid. If you win, your earnings will increase by your value for
the prize, and decrease by the second-highest bid. If you don't
win, your earnings will remain unchanged.
After each auction, we will display all bids. Ties for the
highest bid will be resolved randomly.
\end{lstlisting}
\caption{SPSB auction, base rule (\texttt{private\_second\_price}). Standard second-price sealed-bid auction with independent private values. The FPSB, TPSB, APV, CV, and eBay rule prompts used in the validation exhibits (Section~\ref{sec:validation}) differ only in the allocation/payment paragraph and the value-model paragraphs; the eBay treatment prompts are summarized in Appendix~\ref{app:ebay}.}
\label{fig:prompt-spsb-base}
\end{figure}

\paragraph{Direct DA (traditional description).} The canonical DA baseline (\texttt{da\_direct\_traditional}) describes the full five-step student-proposing algorithm:

\begin{figure}[H]
\begin{lstlisting}
You are Student {{student_id}}. There are 4 schools: w, x, y, z.
You will be matched to at most one school.

HOW THE MATCHING WORKS (Deferred Acceptance):
1. All students simultaneously submit a ranking of schools
2. Each student "proposes" to their top-ranked school
3. Each school tentatively accepts the proposer it ranks highest
   (by priority), rejecting others
4. Rejected students propose to their next choice
5. This repeats until all students are matched

Your preference ordering (most preferred to least preferred):
  {{preference_order}}

Your priority rank at each school (1 = highest priority):
  w: {{pw}}, x: {{px}}, y: {{py}}, z: {{pz}}

Moreover, most applicants seem to favor {{global_ranking}}.

Submit a complete ranking of schools from most preferred to
least preferred.
Format: Ranking: <1st> > <2nd> > <3rd> > <4th>
\end{lstlisting}
\caption{Direct DA, base rule (\texttt{da\_direct\_traditional}). A \emph{null} variant (\texttt{da\_direct\_null}) is identical except that the ``HOW THE MATCHING WORKS'' block is removed entirely, providing no mechanism explanation; its results appear in Appendix~\ref{app:da} (Table~\ref{tab:intervention-inventory-da}).}
\label{fig:prompt-da-traditional}
\end{figure}

\subsection{Family A: extensive-form implementations}\label{app:prompts-osp}

\paragraph{Ascending clock (AC-B), per-tick query.} In the true iterative clock (final ID \path{osp_clock_iterative}), each auction consists of multiple clock cycles; at each posted price every remaining bidder is queried, and no dropout information is displayed, as in \citet{breitmoser2022obviousness}:

\begin{figure}[H]
\begin{lstlisting}
Your value towards to the prize is {value} in this round.
The current price in this clock cycle is {current_price}.
The price for next clock cycle is {current_price + increment}.

Do you want to stay in the bidding?
If you choose yes, you can keep bidding for next clock. If you
choose No, you will exit and have no chance to re-enter the
bidding. Your response must use these EXACT tags below. You
must output the ACTION.
    <PLAN> [...limit 50 words...] </PLAN>
    <ACTION> Yes or No </ACTION>
\end{lstlisting}
\caption{Ascending clock (AC-B), per-tick stay/exit query (final ID \texttt{osp\_clock\_iterative}). The bidder is first shown the auction rules (value model plus clock payment rule); the query above is repeated at every clock tick. See Appendix~\ref{app:procedures} for the full clock simulation loop and the AC (open) variant. On the \texttt{intervention\_proxy\_breitmoser} name collision affecting this lever, see Section~\ref{app:prompts-inventory}.}
\label{fig:prompt-clock-iterative}
\end{figure}

\paragraph{Iterative DA, binary-offer query.} The primary prompt at the binary-offer nodes of the OSP decision tree of \citet{ashlagi2018stable}; it explicitly states the consequence of each answer, providing the ``obviousness witness'' of that construction:

\begin{figure}[H]
\begin{lstlisting}
You are Student {{student_id}}. We are running a step-by-step
matching procedure.

Current remaining schools: {{remaining_set}}

Your preference ordering (most preferred to least preferred):
  {{preference_order}}

Your priority rank at each school (1 = highest priority):
  w: {{pw}}, x: {{px}}, y: {{py}}, z: {{pz}}

Moreover, most applicants seem to favor {{global_ranking}}.

Question: Among the remaining schools ({{remaining_set}}), is
  {{candidate}} your most preferred?
- If you answer YES: you are immediately matched to
  {{candidate}} and leave the process.
- If you answer NO: {{candidate}} remains available to you.
  You continue in the process with all current remaining
  schools ({{fallback_set}}).

<REASON> ... </REASON>
<DECISION> Answer: YES / NO </DECISION>
\end{lstlisting}
\caption{Iterative DA, yes/no query with consequence explanation (\texttt{da\_osp\_yesno}). At serial-dictatorship nodes, a companion prompt (\texttt{da\_osp\_choice}) instead asks the student to name their single most preferred school among the currently available set (``Do not list a full ranking. Do not condition on hypothetical future steps.'').}
\label{fig:prompt-da-osp-yesno}
\end{figure}

\subsection{Family B: mechanism descriptions}\label{app:prompts-desc}

\paragraph{B1 --- Menu restatement (auction).} Based on \citet{gonczarowski2023strategyproofness}, this replaces the payment paragraph of the SPSB base rule with a menu description (\texttt{intervention\_menu}):

\begin{figure}[H]
\begin{lstlisting}
Your "price to win" the item will be set to the highest bid
placed by any other player. If your bid is higher than this
"price to win," then you will win the item and pay this price.
If you don't win, your earnings will remain unchanged.
\end{lstlisting}
\caption{Menu restatement for SPSB (\texttt{intervention\_menu}, Family B1). Reported in Section~\ref{sec:ranking-descriptions} for all four models: small, sign-mixed effects across models and null on average.}
\label{fig:prompt-menu}
\end{figure}

\paragraph{B1 --- Menu restatement (DA).} Two DA analogues exist as templates, \path{da_direct_menu_property} and \path{da_direct_menu_mechanics}; both restate DA through the obtainable-schools menu, following \citet{gonczarowski2023strategyproofness}. Both are reported in Appendix~\ref{app:da}, where they split sharply by content. The \emph{property} variant reads:

\begin{figure}[H]
\begin{lstlisting}
KEY PROPERTY OF THIS MECHANISM:
Given the other students' rankings and school priorities, there
is a fixed set of schools you could possibly be matched with
(your "Obtainable Schools").

IMPORTANT: Your ranking CANNOT change what is in your Obtainable
Schools. Your ranking ONLY determines which school you receive
WITHIN that set - you get your highest-ranked school among the
obtainable ones.
\end{lstlisting}
\caption{Menu restatement for DA, property variant (\texttt{da\_direct\_menu\_property}, Family B1). The \emph{mechanics} variant states the same construction procedurally (``Step 1: \ldots the mechanism computes your Obtainable Schools. Step 2: You receive the school you ranked highest among your Obtainable Schools.'') \emph{without} the invariance statement. Both are reported in Appendix~\ref{app:da}: the property variant significantly improves truthful reporting (pooled $4.2\% \to 1.8\%$), the mechanics variant significantly worsens it ($4.2\% \to 6.9\%$).}
\label{fig:prompt-da-menu}
\end{figure}

\paragraph{B2 --- Payoff Safety (auction).} Surfaces the pivotality property behind \citet{vickrey1961counterspeculation}'s strategy-proofness argument (old labels \path{axis2_forward_onestep} / \path{axis1_contingent_onestep}; final ID \emph{Payoff Safety}):

\begin{figure}[H]
\begin{lstlisting}
**How the auction works:**
Think of your bid as setting your "maximum price." The auction
will find the highest price below your maximum that still wins.
Specifically:
- If your bid is the highest, you win and pay the second-highest
  bid
- If your bid is not the highest, you don't win and pay nothing

This means: Your bid only matters for determining IF you win,
not HOW MUCH you pay. You can think of it simply as: "At what
price would I no longer want the item?"
\end{lstlisting}
\caption{Payoff Safety description for SPSB (Family B2). Despite the ``axis-2 / forward-planning'' repo filename, this is a mechanism description, not a reasoning scaffold: it states \emph{what is true} of the payment rule rather than instructing a mode of reasoning (see the mapping notes in Section~\ref{app:prompts-inventory}).}
\label{fig:prompt-payoff-safety}
\end{figure}

\paragraph{B2 --- Rejection Safety (DA).} Surfaces the property that distinguishes DA from the Boston mechanism \citep{abdulkadirouglu2003school}: rejections are safe (old label \path{axis2_da_monotonic_safety}; final ID \emph{Rejection Safety}):

\begin{figure}[H]
\begin{lstlisting}
KEY INSIGHT - REJECTIONS ONLY REDIRECT, NEVER ELIMINATE:
In this matching process:
- Being rejected from a school you ranked highly does NOT hurt
  your chances at schools ranked lower
- You cannot be "knocked out" of a school that would have
  accepted you
- The algorithm protects you: rejections only happen BEFORE you
  would have been matched

Your ranking determines the ORDER in which you're considered,
not WHETHER you're considered.
\end{lstlisting}
\caption{Rejection Safety description for DA (Family B2). Two sibling ``monotonicity'' templates (\texttt{axis2\_da\_monotonic\_options}: ``your options never shrink''; \texttt{axis2\_da\_monotonic\_outcome}: ``the obtainable set determines the outcome'') were run but are not reported (Table~\ref{tab:intervention-inventory}).}
\label{fig:prompt-rejection-safety}
\end{figure}

\paragraph{B3 --- Clock-framing (auction).} Based on \citet{breitmoser2022obviousness}, this \emph{describes} the sealed bid as an automatic exit threshold in a simulated clock---a description change only; the game remains a one-shot sealed-bid auction (repo template \path{intervention_proxy_breitmoser}; final ID \path{desc_clock_framing}):

\begin{figure}[H]
\begin{lstlisting}
**FIRST STAGE: Sealed Bid**
You will submit a sealed bid privately at the same time as the
other bidders. This bid will serve as your automatic exit price
in the next stage.
**SECOND STAGE: Ascending Clock (Simulation)**
After the sealed bid stage, we will simulate an ascending clock
auction. The clock will start at $0 and increase in increments
of ${{increment}}. The clock will display the current price.
You will also see that there are a total of {{num_bidders}}
bidders participating, although you do not know other bidder's
values. If the current price on the clock reaches or exceeds
your sealed bid from the first stage, you will automatically
exit the auction. The auction ends when only one bidder is left
remaining in the second stage based on their bid from the first
stage.
**END OF AUCTION**
If you win, your earnings will increase by your value for the
prize and decrease by the clock price at the end of the auction.
If you don't win, your earnings will remain unchanged.
\end{lstlisting}
\caption{Clock-framing description for SPSB (Family B3; final ID \texttt{desc\_clock\_framing}). The elicited action is still a single sealed bid; no per-tick queries occur. Never to be conflated with the true iterative clock of Figure~\ref{fig:prompt-clock-iterative}; see the name-collision note in Section~\ref{app:prompts-inventory}.}
\label{fig:prompt-clock-framing}
\end{figure}

\subsection{Family C: reasoning scaffolds}\label{app:prompts-scaffold}

\paragraph{C1 --- Payoff Tree (auction).} Presents the explicit decision tree of outcomes (old labels \path{axis2_forward_tree} and \path{axis1_contingent_tree}; final ID \emph{Payoff Tree}; classified C1 contingent reasoning---see Section~\ref{app:prompts-inventory}):

\begin{figure}[H]
\begin{lstlisting}
**Decision Tree:**
Consider your decision as follows:

You choose bid B. Then one of two things happens:

  PATH A: Your bid B > all other bids
    -> You WIN
    -> You pay: second-highest bid (call it P)
    -> Your profit: (Your Value) - P
    -> Note: P < B, so if Value >= B, then profit >= 0

  PATH B: Your bid B <= some other bid
    -> You LOSE
    -> You pay: $0
    -> Your profit: $0

Key insight: In PATH A, you pay P (not B). So bidding higher
than your value only matters if it changes whether you're in
PATH A vs PATH B.
\end{lstlisting}
\caption{Payoff Tree scaffold for SPSB (Family C1). Non-ASCII arrows in the repo template are transliterated as \texttt{->}.}
\label{fig:prompt-tree-auction}
\end{figure}

\paragraph{C1 --- Payoff Tree (DA).} The DA analogue (\texttt{axis1\_da\_tree}) replaces the algorithm description with a tree view:

\begin{figure}[H]
\begin{lstlisting}
HOW THE MATCHING WORKS (Decision Tree View):
Think of DA as a tree of possibilities:
- ROOT: You propose to your 1st choice
  - If ACCEPTED: You're matched to that school
  - If REJECTED: You propose to your 2nd choice
    - If ACCEPTED: You're matched to that school
    - If REJECTED: You propose to your 3rd choice
      - And so on...

At each branch, acceptance/rejection depends on the school's
priorities and who else proposed there.
\end{lstlisting}
\caption{Payoff Tree scaffold for DA (Family C1). Additional C1 variants in both domains (\emph{Enumerate}, \emph{Dominated}, \emph{Worst-Case}, \emph{Backward-Induction} framings) were run but are not reported; see Table~\ref{tab:intervention-inventory}. The \emph{Enumerate} variant, for example, asks: ``What are the possible bids the other players might submit? For each possible bid they might make, what would be your best response?'' (auctions), or the analogous question over rankings (DA).}
\label{fig:prompt-tree-da}
\end{figure}

\paragraph{C2 --- Forward-planning lookahead (DA only).} Scaffolds $k$-step simulation of the DA algorithm's internal rounds \citep{pycia2023theory}; SPSB has no analogous internal dynamics, so C2 is DA-only (Section~\ref{sec:framework}). The $k{=}1$ variant (\texttt{axis2\_da\_1step}):

\begin{figure}[H]
\begin{lstlisting}
HOW THE MATCHING WORKS:
You submit a ranking. The algorithm processes your ranking
sequentially:
- First, you "propose" to your top-ranked school
- If accepted, you're matched there
- If rejected, you propose to your second choice

THINK ONE STEP AHEAD:
Before finalizing your ranking, consider: If you are rejected
from your first choice, which school would you propose to next?
Does this affect how you should rank schools?
\end{lstlisting}
\caption{One-step lookahead scaffold for DA (Family C2, \texttt{axis2\_da\_1step}). The $k{=}2$ variant (\texttt{axis2\_da\_2step}) extends the framing to the first two potential rejections (``Plan your ranking with these two steps in mind''). The $k{=}0$ and $k{=}\infty$ (full-simulation) variants were run but are not reported.}
\label{fig:prompt-lookahead}
\end{figure}

\paragraph{C3 --- Belief scaffolds (both domains).} First- and second-order belief prompts appended to the base rules. In auctions (\texttt{axis3\_beliefs\_firstorder}, \texttt{axis3\_beliefs\_secondorder}):

\begin{figure}[H]
\begin{lstlisting}
[First-order]
The other bidders are rational agents who are also trying to
maximize their earnings.
Before bidding, ask yourself: What do I expect the other bidders
to bid? Given their values are also drawn randomly between $0
and ${{private}}, and they want to maximize their earnings, how
do I think they will behave? Does this affect what I should bid?

[Second-order]
Before bidding, consider the following: What do the other
bidders think YOU will bid? They know you are rational and
trying to maximize your earnings. They might try to anticipate
your strategy. Does their belief about your behavior affect what
they will do? And does that, in turn, affect what you should do?
\end{lstlisting}
\caption{Belief scaffolds for SPSB (Family C3). The DA analogues (\texttt{axis3\_da\_firstorder}, \texttt{axis3\_da\_secondorder}) ask ``What rankings do you think the other students will submit?'' and ``What do you think the other students believe YOU will rank?'' respectively. A \emph{Common Knowledge} variant in each domain (emphasizing common knowledge of rationality) was run but is not reported.}
\label{fig:prompt-beliefs}
\end{figure}

\subsection{Family D: preference framings}\label{app:prompts-prospect}

The Family-D battery (results in Appendix~\ref{app:prospect}) replaces or augments the objective framing rather than the mechanism explanation. As one representative, the risk-averse persona (\path{intervention_risk_averse}, identical wording in both domains modulo context):

\begin{figure}[H]
\begin{lstlisting}
You are risk-averse. This means that you are the type of person
that would only pay $4 for a coin toss where you got $0.00 on
tails and $10.00 on heads (even though the expected value is
$5.00). You prefer to avoid uncertainty and would rather have a
guaranteed outcome than take a risky gamble, even if the gamble
has the same or slightly higher expected value.
\end{lstlisting}
\caption{Risk-averse persona (Family D). Risk-neutral and risk-seeking personas change only the coin-toss price (\$5, \$6); the loss-aversion battery (gain/loss/mixed frames, endowment, WTA/WTP) is described in Appendix~\ref{app:prospect-loss}.}
\label{fig:prompt-risk-persona}
\end{figure}

\subsection{Intervention inventory and old-to-new mapping}\label{app:prompts-inventory}

The prompt templates accumulated across two experimental programs and several repo versions, and three of the resulting labels are actively misleading. We fix the mapping here once, and Table~\ref{tab:intervention-inventory} applies it to every template in the repo.

\paragraph{Taxonomy corrections (plan reconciliation R7).} (i)~\emph{Payoff Tree} carries the repo filename \path{axis2_forward_tree.txt} (V12) / \path{axis1_contingent_tree.txt} (imported pipeline), but SPSB is a one-shot game with no own future decision points, so there is no forward-planning axis to scaffold; the intervention lays out payoff-relevant \emph{contingencies} and is classified C1. (ii)~\emph{Payoff Safety} (\path{axis2_forward_onestep} / \path{axis1_contingent_onestep}) and \emph{Rejection Safety} (\path{axis2_da_monotonic_safety}) state true properties of the payment/rejection rule without instructing any mode of reasoning; both are descriptions (B2), not scaffolds. (iii)~Clock-framing is a description (B3), not an extensive-form change.

\paragraph{The \texttt{intervention\_proxy\_breitmoser} name collision (plan reconciliation R6).} The single template name \path{intervention_proxy_breitmoser} has been used for two different objects: (a)~in the ES appendix, it labels the \emph{true iterative clock} prompt (per-tick stay/exit queries, Figure~\ref{fig:prompt-clock-iterative}); (b)~in the repo template files (\path{rule_template/V10/}, \path{V12/}, and the imported pipeline's \path{rule_template/auctions/}), it is the \emph{two-stage sealed-bid clock-framing description} (Figure~\ref{fig:prompt-clock-framing}). We split the ID: \path{osp_clock_iterative} (Family A) vs.\ \path{desc_clock_framing} (Family B3). Our audit of the imported pipeline supports the split as follows: the logs stored under \path{experiment_logs/<model>/intervention_proxy_breitmoser/} record exactly one sealed bid per bidder per auction (no tick-by-tick decisions), i.e., they are \path{desc_clock_framing} runs; whereas the clock column of the OSP comparison (Figure~\ref{fig:osp-comparison}) reproduces numerically from the pipeline's \path{ascending_clock_closed} cells, which are true per-tick clock runs (in the APV environment---see the provenance flag in Section~\ref{sec:ranking-extensive}). The audit is settled: the pipeline template behind \path{intervention_proxy_breitmoser} is the \emph{description} variant (\path{desc_clock_framing}), and the OSP clock column is computed from the \path{ascending_clock_closed} cells. With both IDs now reported for all four models (Sections~\ref{sec:ranking-extensive} and~\ref{sec:ranking-descriptions}), ranking rungs A (true clock) and B3 (clock-framing) are empirically distinguished model by model.

\paragraph{Run-status audit.} Table~\ref{tab:intervention-inventory} records, for every declared template, whether its results are reported in this paper (with the section), were \emph{run but not reported} (logs exist in \path{engineer_simplicity-main/experiment_logs/} or in the git-recovered writeup logs, but the cells do not appear in any exhibit), or were \emph{not run} on the canonical grid. This implements the no-phantom-results rule (plan \S7.8): no template is silently declared.

\paragraph{Comparison-baseline declaration.} Two families of sealed-bid SPSB baselines coexist in the pipeline, and the paper uses them for different purposes: OSP and extensive-form comparisons quote the \emph{dedicated} \texttt{spsb} cells (mean deviations $-0.49$/$-1.63$/$-3.28$/$-5.27$ for Claude/Gemini/GPT-4o/Gemma, matching the previously published $-0.5$/$-1.6$/$-3.3$/$-5.3$), while intervention contrasts use the \emph{axis-specific} baselines (each \texttt{axisK\_*} treatment against its same-axis baseline; the pooled axis baseline for the \texttt{intervention\_*} cells). Both baseline sets, and every treatment cell, are recorded with bootstrap confidence intervals in \path{results/merged_ranking/auction_cells.csv}.

\begin{table}[p]
\centering
\footnotesize
\setlength{\tabcolsep}{4pt}
\begin{tabular}{p{4.9cm} p{3.3cm} p{1.0cm} p{4.4cm}}
\toprule
\textbf{Old label (repo template)} & \textbf{Final taxonomy family} & \textbf{Domain} & \textbf{Reported in paper?} \\
\midrule
\multicolumn{4}{l}{\emph{Panel A: Auctions}} \\
\midrule
\texttt{private\_second\_price} & Baseline (SPSB, IPV) & Auction & Yes (\S\ref{sec:validation}, \S\ref{sec:ranking}) \\
\texttt{affiliated\_spsb}; \texttt{affiliated\_ac}; \texttt{affiliated\_ac\_closed} & Baseline (APV); A (AC); A (AC-B) & Auction & Yes (\S\ref{sec:validation}, \S\ref{sec:ranking-extensive}) \\
\texttt{intervention\_proxy\_breitmoser} (ES-appendix usage) $\to$ \texttt{osp\_clock\_iterative} & A (extensive form) & Auction & Yes (\S\ref{sec:ranking-extensive}); clock column = \texttt{ascending\_clock\_closed} cells (provenance settled; see below) \\
\texttt{intervention\_proxy\_breitmoser} (repo template) $\to$ \texttt{desc\_clock\_framing} & B3 (clock-framing) & Auction & Yes, all four models (\S\ref{sec:ranking-descriptions}); large improvement in every model ($p<0.001$) \\
\texttt{intervention\_menu} & B1 (menu restatement) & Auction & Yes, all four models (\S\ref{sec:ranking-descriptions}); sign-mixed, null on average \\
\texttt{axis2\_forward\_onestep} = \texttt{axis1\_contingent\_onestep} $\to$ Payoff Safety & B2 (safety-exposing) & Auction & Yes (\S\ref{sec:ranking-descriptions}) \\
\texttt{axis2\_forward\_tree} = \texttt{axis1\_contingent\_tree} $\to$ Payoff Tree & C1 (contingent) & Auction & Yes (\S\ref{sec:ranking-scaffolds}) \\
\texttt{axis1\_contingent\_\{enumerate, dominated, worstcase\}} & C1 (contingent) & Auction & Run; the \emph{Worst-Case} cell enters the trace analysis as a manipulation check (App.~\ref{app:traces}); \emph{Enumerate}/\emph{Dominated} not reported \\
\texttt{axis2\_forward\_backward\_induct} = \texttt{axis1\_contingent\_backward\_induct} & C1 (two-stage framing) & Auction & Run, not reported \\
\texttt{axis3\_beliefs\_firstorder}; \texttt{axis3\_beliefs\_secondorder} & C3 (beliefs) & Auction & Yes (\S\ref{sec:ranking-scaffolds}) \\
\texttt{axis3\_beliefs\_common\_knowledge} & C3 (beliefs) & Auction & Run, not reported \\
\texttt{axis1\_contingent\_baseline}; \texttt{axis2\_forward\_baseline}; \texttt{axis3\_beliefs\_baseline}; \texttt{loss\_aversion\_baseline} & Per-axis baselines & Auction & Yes, as comparison baselines for axis-matched treatment cells (see the comparison-baseline declaration above) \\
\texttt{intervention\_risk\_\{averse, neutral, seeking\}} & D (risk personas) & Auction & Yes (App.~\ref{app:prospect}) \\
\texttt{loss\_aversion\_\{gain\_frame, loss\_frame, mixed\_frame, endowment, WTA\_WTP\}} & D (loss/endowment frames) & Auction & Yes (App.~\ref{app:prospect}) \\
\texttt{intervention\_NE\_strat\_reveal}; \texttt{intervention\_wrong\_strat\_reveal}; \texttt{intervention\_nash\_deviation} & Excluded: advice (reveals or misstates the strategy; not a simplicity-preserving lever) & Auction & Run, not reported \\
\texttt{first\_price\_interventions/*}; \texttt{third\_price\_interventions/*} & Out of ranking scope (BNE benchmark, not dominant-strategy) & Auction & Partially run, not reported; parked per plan \S6 scope decision \\
\texttt{robust\_*} (currency/language); \texttt{private\_all\_pay*} & Legacy V10 robustness & Auction & Run in earlier drafts, not reported (cut, plan \S7.9) \\
\bottomrule
\end{tabular}
\caption{Intervention inventory (Panel A: auctions): every declared prompt template, its final taxonomy family (Section~\ref{sec:framework}), and its run/reporting status. ``Run, not reported'' means raw logs exist in the imported pipeline (\texttt{Engineering\_simplicity/engineer\_simplicity-main/experiment\_logs/}) or the git-recovered writeup logs, but the cell appears in no exhibit of this paper. Old labels are repo filenames; where the V12 writeup name and the imported-pipeline name differ, both are given with ``=''. Panel B (deferred acceptance) continues in Table~\ref{tab:intervention-inventory-da}.}
\label{tab:intervention-inventory}
\end{table}

\begin{table}[p]
\centering
\footnotesize
\setlength{\tabcolsep}{4pt}
\begin{tabular}{p{4.9cm} p{3.3cm} p{1.0cm} p{4.4cm}}
\toprule
\textbf{Old label (repo template)} & \textbf{Final taxonomy family} & \textbf{Domain} & \textbf{Reported in paper?} \\
\midrule
\multicolumn{4}{l}{\emph{Panel B: Deferred acceptance}} \\
\midrule
\texttt{da\_direct\_traditional} & Baseline (direct DA) & DA & Yes (App.~\ref{app:da}) \\
\texttt{da\_direct\_null} & Baseline variant (no explanation) & DA & Yes (App.~\ref{app:da}); heterogeneous across models \\
\texttt{da\_direct\_textbook\_sp} & Direct SP statement (reveals the dominant strategy) & DA & Yes, as a diagnostic ceiling (App.~\ref{app:da}): $0.0\%$ error for all four models; excluded from the lever ranking \\
\texttt{da\_direct\_menu\_property}; \texttt{da\_direct\_menu\_mechanics} & B1 (menu restatement) & DA & Yes (App.~\ref{app:da}); property variant improves play, mechanics variant backfires \\
\texttt{da\_osp\_yesno} (+ \texttt{da\_osp\_choice} node prompt; \texttt{da\_osp\_yesno\_guaranteed}) & A (iterative DA) & DA & Yes (App.~\ref{app:da}); the pick-based protocol is the canonical iterative-DA cell, the yes/no tree appears in the decision-level accounting \\
\texttt{axis1\_da\_tree} $\to$ Payoff Tree (DA) & C1 (contingent) & DA & Yes (App.~\ref{app:da}) \\
\texttt{axis1\_da\_\{enumerate, dominated, worstcase, onestep, backward\_induct\}} & C1 (contingent) & DA & Run, not reported \\
\texttt{axis2\_da\_1step}; \texttt{axis2\_da\_2step} & C2 (forward planning) & DA & Yes (App.~\ref{app:da}) \\
\texttt{axis2\_da\_0step}; \texttt{axis2\_da\_fullsim} & C2 (forward planning) & DA & Run, not reported \\
\texttt{axis2\_da\_monotonic\_safety} $\to$ Rejection Safety & B2 (safety-exposing) & DA & Yes (App.~\ref{app:da}) \\
\texttt{axis2\_da\_monotonic\_options}; \texttt{axis2\_da\_monotonic\_outcome} & B2 variants & DA & Run, not reported \\
\texttt{axis3\_da\_firstorder}; \texttt{axis3\_da\_secondorder} & C3 (beliefs) & DA & Yes (App.~\ref{app:da}) \\
\texttt{axis3\_da\_common\_knowledge} & C3 (beliefs) & DA & Run, not reported \\
\texttt{intervention\_risk\_\{averse, neutral, seeking\}} (DA) & D (risk personas) & DA & Yes (App.~\ref{app:prospect}) \\
\texttt{loss\_aversion\_*} (DA) & D (loss/endowment frames) & DA & Yes (App.~\ref{app:prospect}) \\
\texttt{da\_cardinal} variants (pipeline) & Environment variant (cardinal payoffs) & DA & Run, not reported \\
\bottomrule
\end{tabular}
\caption{Intervention inventory, continued (Panel B: deferred acceptance). Conventions as in Table~\ref{tab:intervention-inventory}. The eBay treatment prompts (T1--T4) belong to the eBay environment rather than the lever taxonomy and are catalogued in Appendix~\ref{app:ebay}.}
\label{tab:intervention-inventory-da}
\end{table}
